\documentclass[11pt]{article}

\usepackage[margin=1in]{geometry}
\usepackage{amsmath, amssymb, amsthm}
\usepackage{graphicx}
\usepackage{hyperref}
\usepackage{enumitem}
\usepackage{xcolor}
\usepackage{listings}
\usepackage{booktabs}
\usepackage{parskip}

% allow LaTeX extra flexibility so long \texttt{...} runs don't overflow the margin
\setlength{\emergencystretch}{3em}
\sloppy

% figure placeholder: shows a framed box until the real PNG is added to figures/,
% then automatically renders the image. Drop the exported W&B chart at the given path.
\newcommand{\figph}[1]{%
  \IfFileExists{#1}%
    {\includegraphics[width=0.8\linewidth]{#1}}%
    {\fbox{\parbox[c][3.5cm][c]{0.8\linewidth}{\centering\itshape Figure placeholder --- export W\&B curve to \texttt{\detokenize{#1}}}}}%
}

\lstset{
  basicstyle=\ttfamily\small,
  breaklines=true,
  frame=single,
  columns=fullflexible,
}

\newenvironment{answer}%
  {\par\noindent\textbf{Answer.}\par\vspace{0.25em}\begingroup\color{black!85}}%
  {\endgroup\par\vspace{1em}}

\newcommand{\problem}[2]{%
  \subsection*{Problem (\texttt{#1}) --- #2}%
  \addcontentsline{toc}{subsection}{#1 --- #2}%
}

\title{CS336 Assignment 1: Basics \\ Written Deliverables}
\author{Rohan Rajput}
\date{\today}

\begin{document}
\maketitle
\tableofcontents
\newpage


\section{Byte-Pair Encoding (BPE) Tokenizer}

\problem{unicode1}{Understanding Unicode (1 point)}

\textbf{(a)} What Unicode character does \texttt{chr(0)} return? \\
\emph{Deliverable: A one-sentence response.}
\begin{answer}
\verb|'\x00'|
\end{answer}

\textbf{(b)} How does this character's string representation (\texttt{\_\_repr\_\_()}) differ from its printed representation? \\
\emph{Deliverable: A one-sentence response.}
\begin{answer}
 \texttt{\_\_repr\_\_()} adds quotes and escape sequences for non-printable characters, while
  \texttt{print} emits raw characters.
\end{answer}

\textbf{(c)} What happens when this character occurs in text? (Try \texttt{chr(0)}, \texttt{print(chr(0))}, and the concatenation examples in the handout.) \\
\emph{Deliverable: A one-sentence response.}
\begin{answer}
\textit{It adds black space in the text}
\end{answer}


\problem{unicode2}{Unicode Encodings (3 points)}

\textbf{(a)} What are some reasons to prefer training our tokenizer on UTF-8 encoded bytes, rather than UTF-16 or UTF-32? \\
\emph{Deliverable: A one-to-two sentence response.}
\begin{answer}
\textit{UTF-8 is preferred for BPE tokenization because it's variable-width (compact for common ASCII text, reducing sequence lengths and merges needed), self-synchronizing (high-bit prefixes mark character boundaries, enabling corruption recovery and mid-stream parsing), and endianness-free (no BOM required, unlike UTF-16/32).}
\end{answer}

\textbf{(b)} The function \verb|decode_utf8_bytes_to_str_wrong| in the handout is incorrect. Why? Provide an example input byte string that yields incorrect results. \\
\emph{Deliverable: An example input byte string for which the function produces incorrect output, with a one-sentence explanation of why it is incorrect.}
\begin{answer}
\textit{"hello é, it gives error because UTF-8 is a variable length encoding and non-ascii character bytes can't be decode indivisually"}
\end{answer}

\textbf{(c)} Give a two-byte sequence that does not decode to any Unicode character(s). \\
\emph{Deliverable: An example, with a one-sentence explanation.}
\begin{answer}
\textit{""}
\end{answer}


\problem{train\_bpe}{BPE Tokenizer Training (15 points)}

Write a function that, given a path to an input text file, trains a (byte-level) BPE tokenizer with the interface described in the handout (inputs: \texttt{input\_path}, \texttt{vocab\_size}, \texttt{special\_tokens}; outputs: \texttt{vocab}, \texttt{merges}). Implement \texttt{[adapters.run\_train\_bpe]} and pass \texttt{uv run pytest tests/test\_train\_bpe.py}.

\emph{Deliverable: Code that passes the provided tests. (No written answer required --- include any implementation notes below if desired.)}
\begin{answer}
% TODO (optional notes)
\end{answer}


\problem{train\_bpe\_tinystories}{BPE Training on TinyStories (2 points)}

\textbf{(a)} Train a byte-level BPE tokenizer on TinyStories (vocab size 10{,}000, with \texttt{<|endoftext|>} as a special token). How much time and memory did training take? What is the longest token in the vocabulary --- does it make sense? \\
\emph{Deliverable: A one-to-two sentence response.}
\begin{answer}
- Time taken: 43.7s
- Peak Memory: 0.18GB
- Longest token : b' accomplishment'
\end{answer}

\textbf{(b)} Profile your code. What part of the tokenizer training process takes the most time? \\
\emph{Deliverable: A one-to-two sentence response.}
\begin{answer}
Pre-token : 19s wall, 200s total cpu across 12 workers
Merge loop: 25s wall
\end{answer}


\problem{train\_bpe\_expts\_owt}{BPE Training on OpenWebText (2 points)}

\textbf{(a)} Train a byte-level BPE tokenizer on OpenWebText with vocab size 32{,}000. What is the longest token in the vocabulary --- does it make sense? \\
\emph{Deliverable: A one-to-two sentence response.}
\begin{answer}
The longest token is 64 length long as mojibake whitespace scraping artifact
\end{answer}

\textbf{(b)} Compare and contrast the tokenizer trained on TinyStories vs. OpenWebText. \\
\emph{Deliverable: A one-to-two sentence response.}
\begin{answer}
Tiny story: 15 bytes 'accomplishment', a real meaningful word
Owt : 64 bytes a mojibake whitespace scrapping effect 
\end{answer}


\problem{tokenizer}{Implementing the tokenizer (15 points)}

Implement a \texttt{Tokenizer} class supporting \texttt{\_\_init\_\_}, \texttt{from\_files}, \texttt{encode}, \texttt{encode\_iterable}, and \texttt{decode}, with special-token support. Implement \texttt{[adapters.get\_tokenizer]} and pass \texttt{uv run pytest tests/test\_tokenizer.py}.

\emph{Deliverable: Code that passes the provided tests.}
\begin{answer}
% TODO (optional notes)
\end{answer}


\problem{tokenizer\_experiments}{Experiments with tokenizers (4 points)}

\textbf{(a)} Sample 10 documents from TinyStories and OpenWebText. What is each tokenizer's compression ratio (bytes/token)? \\
\emph{Deliverable: A one-to-two sentence response.}
\begin{answer}
On a sample of 10 validation documents, the TinyStories 10k tokenizer achieves $\approx 4.05$ bytes/token, and the OpenWebText 32k tokenizer achieves $\approx 4.25$ bytes/token. The larger OWT vocabulary gives only a marginally better ratio despite being $3\times$ bigger, because OWT text is far more diverse, so the extra vocabulary is spent covering variety rather than tightly compressing any single domain.
\end{answer}

\textbf{(b)} What happens if you tokenize your OpenWebText sample with the TinyStories tokenizer? Compare compression ratio / describe qualitatively. \\
\emph{Deliverable: A one-to-two sentence response.}
\begin{answer}
The compression ratio goes down (3.16) because tinystories merges don't cover the OWT vocab, so text fragments towards byte-level 
tokens i.e. more tokens per bytes
\end{answer}

\textbf{(c)} Estimate the throughput of your tokenizer (bytes/second). How long to tokenize the Pile (825\,GB)? \\
\emph{Deliverable: A one-to-two sentence response.}
\begin{answer}
Our throughput is 1.04 MB/s so it will take approx 9.9 days for the Pile.
\end{answer}

\textbf{(d)} Encode the TinyStories and OpenWebText train/dev sets to integer token IDs (serialize as NumPy \texttt{uint16}). Why is \texttt{uint16} an appropriate choice? \\
\emph{Deliverable: A one-to-two sentence response.}
\begin{answer}
The numpy uint16 takes less space compare to uint32/64. It can fit entire vocab in 0-65535 digit.
\end{answer}


\section{Transformer Language Model Architecture}

\problem{linear}{Implementing the linear module (1 point)}
Implement a \texttt{Linear} class (subclass of \texttt{torch.nn.Module}) following the interface in the handout. Implement the adapter and pass the provided tests.
\emph{Deliverable: Code that passes the provided tests.}
\begin{answer}
% TODO (optional notes)
\end{answer}


\problem{embedding}{Implement the embedding module (1 point)}
Implement an \texttt{Embedding} class (subclass of \texttt{torch.nn.Module}).
\emph{Deliverable: Code that passes the provided tests.}
\begin{answer}
% TODO (optional notes)
\end{answer}


\problem{rmsnorm}{Root Mean Square Layer Normalization (1 point)}
Implement RMSNorm as a \texttt{torch.nn.Module}.
\emph{Deliverable: Code that passes the provided tests.}
\begin{answer}
% TODO (optional notes)
\end{answer}


\problem{positionwise\_feedforward}{Implement the position-wise feed-forward network (2 points)}
Implement the SwiGLU feed-forward network (SiLU activation + GLU).
\emph{Deliverable: Code that passes the provided tests.}
\begin{answer}
% TODO (optional notes)
\end{answer}


\problem{rope}{Implement RoPE (2 points)}
Implement a \texttt{RotaryPositionalEmbedding} class.
\emph{Deliverable: Code that passes the provided tests.}
\begin{answer}
% TODO (optional notes)
\end{answer}


\problem{softmax}{Implement softmax (1 point)}
Write a function that applies softmax along a given dimension (numerically stable).
\emph{Deliverable: Code that passes the provided tests.}
\begin{answer}
% TODO (optional notes)
\end{answer}


\problem{scaled\_dot\_product\_attention}{Implement scaled dot-product attention (5 points)}
Implement scaled dot-product attention with optional boolean mask, supporting arbitrary batch-like dimensions.
\emph{Deliverable: Code that passes the provided tests.}
\begin{answer}
% TODO (optional notes)
\end{answer}


\problem{multihead\_self\_attention}{Implement causal multi-head self-attention (5 points)}
Implement causal multi-head self-attention as a \texttt{torch.nn.Module}, with parameters \texttt{d\_model}, \texttt{num\_heads}, etc.
\emph{Deliverable: Code that passes the provided tests.}
\begin{answer}
% TODO (optional notes)
\end{answer}


\problem{transformer\_block}{Implement the Transformer block (3 points)}
Implement the pre-norm Transformer block.
\emph{Deliverable: Transformer block code that passes the provided tests.}
\begin{answer}
% TODO (optional notes)
\end{answer}


\problem{transformer\_lm}{Implementing the Transformer LM (3 points)}
Implement the full Transformer language model (embedding $\rightarrow$ \texttt{num\_layers} blocks $\rightarrow$ final norm $\rightarrow$ LM head).
\emph{Deliverable: A Transformer LM module that passes the provided tests.}
\begin{answer}
% TODO (optional notes)
\end{answer}


\problem{transformer\_accounting}{Transformer LM resource accounting (5 points)}

\textbf{(a)} For a GPT-2 XL-shaped model (\texttt{vocab\_size}=50{,}257, \texttt{context\_length}=1{,}024, \texttt{num\_layers}=48, \texttt{d\_model}=1{,}600, \texttt{num\_heads}=25, \texttt{d\_ff}=4{,}288): how many trainable parameters, and how much memory to load (fp32)? \\
\emph{Deliverable: A one-to-two sentence response.}
\begin{answer}
Summing all parameters (token embedding $Vd$; per block: $2d$ RMSNorm gains, $4d^2$ attention QKVO, $3\,d\,d_{\text{ff}}$ SwiGLU; final RMSNorm $d$; LM head $Vd$, with untied embeddings) gives
\[
P \approx 1.64\text{B trainable parameters}.
\]
At $4$ bytes/parameter (fp32), loading the model requires $4P \approx 6.56$ GB.
\end{answer}

\textbf{(b)} List the matrix multiplies required for one forward pass of GPT-2 XL. Total FLOPs? \\
\emph{Deliverable: A list of matrix multiplies (with descriptions), and the total number of FLOPs required.}
\begin{answer}
Using the convention that a matmul $(m\times n)\cdot(n\times p)$ costs $2mnp$ FLOPs, with $\text{seq}=1024$ (element-wise ops such as norms, SiLU, and softmax are negligible):

\begin{center}
\begin{tabular}{lll}
\toprule
Matmul & Formula & FLOPs \\
\midrule
Q, K, V projections ($\times 3$) & $3\cdot 2\,\text{seq}\,d^2$ & $15.73$ B \\
Attention scores $QK^\top$ & $2\,\text{seq}^2 d$ & $3.36$ B \\
Attention $\times V$ & $2\,\text{seq}^2 d$ & $3.36$ B \\
Output projection & $2\,\text{seq}\,d^2$ & $5.24$ B \\
SwiGLU FFN ($W_1,W_2,W_3$) & $6\,\text{seq}\,d\,d_{\text{ff}}$ & $42.15$ B \\
\midrule
Per block & & $69.84$ B \\
$\times\, 48$ layers & & $3.352$ T \\
LM head & $2\,\text{seq}\,d\,V$ & $164.7$ B \\
\bottomrule
\end{tabular}
\end{center}

Total forward-pass cost $\approx 3.52$ TFLOPs (consistent with the $2 \cdot P \cdot \text{tokens}$ rule of thumb: $2 \times 1.64\text{B} \times 1024 \approx 3.4$ TFLOPs).
\end{answer}

\textbf{(c)} Which parts of the model require the most FLOPs? \\
\emph{Deliverable: A one-to-two sentence response.}
\begin{answer}
The SwiGLU feed-forward network dominates, at $\approx 60\%$ of each block's FLOPs ($42.15$ B of $69.84$ B). The linear projections (QKV + output) are next at $\approx 30\%$, and the actual attention $\text{seq}^2$ matmuls (scores and $\times V$) are smallest at $\approx 10\%$. At $\text{seq}=1024$ the model is matmul-bound on the position-wise feed-forward and projections, not on attention.
\end{answer}

\textbf{(d)} Repeat for GPT-2 small (12L, 768, 12h), medium (24L, 1024, 16h), and large (36L, 1280, 20h). As model size grows, which components take proportionally more/less of total FLOPs? \\
\emph{Deliverable: For each model, a breakdown of components and their FLOPs (as a proportion of total), plus a one-to-two sentence description.}
\begin{answer}
Per-block FLOP share for each model (with $d_{\text{ff}} = \tfrac{8}{3}d$ rounded to a multiple of 64: small $2048$, medium $2752$, large $3392$, XL $4288$):

\begin{center}
\begin{tabular}{lcccc}
\toprule
Model & $d_{\text{model}}$ & FFN & Attn.\ proj.\ (QKVO) & Attn.\ scores \\
\midrule
small  & 768  & $54.5\%$ & $27.3\%$ & $18.2\%$ \\
medium & 1024 & $57.3\%$ & $28.4\%$ & $14.2\%$ \\
large  & 1280 & $58.7\%$ & $29.5\%$ & $11.8\%$ \\
XL     & 1600 & $60.3\%$ & $30.0\%$ & $9.6\%$  \\
\bottomrule
\end{tabular}
\end{center}

As $d_{\text{model}}$ grows (with $\text{seq}$ fixed), the feed-forward and projection matmuls take a larger proportion because they scale as $d_{\text{model}}^2$, whereas the attention score/value matmuls scale as $\text{seq}^2 \cdot d_{\text{model}}$ (only linear in $d_{\text{model}}$) and therefore shrink proportionally.
\end{answer}

\textbf{(e)} Take GPT-2 XL and increase context length to 16{,}384. How does total forward-pass FLOPs change? Relative contributions? \\
\emph{Deliverable: A one-to-two sentence response.}
\begin{answer}
Increasing $\text{seq}$ from $1024$ to $16{,}384$ ($16\times$) raises the total forward-pass cost from $3.52$ TFLOPs to $\approx 133.6$ TFLOPs --- about $38\times$, more than $16\times$, because the attention score/value terms scale as $\text{seq}^2$ ($\times 256$) while the projections/FFN/LM head scale only linearly in $\text{seq}$ ($\times 16$). Consequently the attention score/value share balloons from $\approx 10\%$ to $\approx 63\%$ of the FLOPs, making attention the dominant cost at long context --- the reverse of the trend in part (d).
\end{answer}


\section{Training a Transformer LM}

\problem{cross\_entropy}{Implement cross-entropy (1 point)}
Write a numerically stable cross-entropy function; handle batch dimensions.
\emph{Deliverable: Code that passes the provided tests.}
\begin{answer}
% TODO (optional notes)
\end{answer}


\problem{learning\_rate\_tuning}{Tuning the learning rate (1 point)}
Run the SGD toy example with learning rates 1e1, 1e2, 1e3 for 10 iterations. What happens?
\emph{Deliverable: A one-to-two sentence response with the behaviors you observed.}
\begin{answer}
The learning rate controls how fast the loss converges, up to a stability limit. At $\text{lr}=10$ the loss decays steadily but slowly (from $\approx 26.3$ to $\approx 3.5$ over 10 iterations, not yet converged); at $\text{lr}=100$ it converges fastest, collapsing to near machine zero ($\approx 10^{-16}$) within about 4 iterations; and at $\text{lr}=1000$ it diverges, with the loss exploding without bound ($26 \to 9.5\times10^3 \to 1.6\times10^6 \to \dots \to 2.4\times10^{18}$) because each step overshoots the minimum and the updates amplify.
\end{answer}


\problem{adamw}{Implement AdamW (2 points)}
Implement AdamW as a subclass of \texttt{torch.optim.Optimizer}.
\emph{Deliverable: Code that passes the provided tests.}
\begin{answer}
% TODO (optional notes)
\end{answer}


\problem{adamw\_accounting}{Resource accounting for training with AdamW (2 points)}

\textbf{(a)} Peak memory (float32): decompose into parameters, activations, gradients, optimizer state. Express in terms of \texttt{batch\_size} and model hyperparameters (assume $d_{\text{ff}} = \tfrac{8}{3} d_{\text{model}}$). \\
\emph{Deliverable: An algebraic expression for each of parameters, activations, gradients, and optimizer state, as well as the total.}

\begin{answer}

\textbf{Parameter, gradient, and optimizer memory (bytes):}
\begin{align*}
\text{Parameters} &= 4P = 4 \times 1.6\text{B} \\
\text{Gradients} &= 4P = 4 \times 1.6\text{B} \\
\text{Optimizer state} &= 2 \cdot 4P = 8P \\
\text{Total} &= 4P + 4P + 8P = 16P
\end{align*}

\textbf{Activation memory (per layer):}
\begin{align*}
\text{RMSNorm} &= 2 \cdot B \cdot L \cdot d_{\text{model}} \\
\text{QKV projections} &= 3 \cdot B \cdot L \cdot d_{\text{model}} \\
\text{Attention scores + softmax} &= 2 \cdot B \cdot n_{\text{heads}} \cdot L^2 \\
\text{Attn} \times V \text{ + output proj} &= 2 \cdot B \cdot L \cdot d_{\text{model}} \\
\text{SwiGLU} &= 4 \cdot B \cdot L \cdot d_{\text{ff}} + B \cdot L \cdot d_{\text{model}}
\end{align*}

\textbf{Total memory:}
\begin{align*}
\text{Total} &= 16P \\
&\quad + 4 \Big[ n_{\text{layers}} \left( 8\, B L\, d_{\text{model}} + 2\, B\, n_{\text{heads}} L^2 + 4\, B L\, d_{\text{ff}} \right) \\
&\qquad\quad + B L\, d_{\text{model}} + 2\, B L\, \text{vocab\_size} \Big] \quad \text{bytes}
\end{align*}

\end{answer}

\textbf{(b)} Instantiate for GPT-2 XL; find max batch size that fits in 80\,GB. \\
\emph{Deliverable: An expression of the form $a \cdot \text{batch\_size} + b$ for numerical $a, b$, and the max batch size.}
\begin{answer}
For GPT-2 XL ($P=1.64\text{B}$, $d_{\text{model}}=1600$, $L=1024$, $n_{\text{heads}}=25$, $n_{\text{layers}}=48$, $\text{vocab}=50257$, $d_{\text{ff}}=4288$), the total memory is
\[
\text{Total} = 16.37 \cdot \text{batch\_size} + 26.24 \ \text{GB},
\]
where $b = 16P = 26.24$ GB (parameters + gradients + optimizer state) and $a = 16.37$ GB is the per-batch activation cost.

Solving for the 80 GB budget:
\[
16.37 \cdot \text{batch\_size} + 26.24 \le 80 \;\Longrightarrow\; \text{batch\_size} \le 3.28.
\]
Since batch size must be a positive integer, the maximum is $\boxed{3}$.
\end{answer}

\textbf{(c)} FLOPs per AdamW step? \\
\emph{Deliverable: An algebraic expression, with brief justification.}
\begin{answer}
Let $F$ denote the FLOPs of one forward pass ($\approx 3.52$ TFLOPs for GPT-2 XL). The backward pass costs $2F$, and the AdamW parameter update is a handful of element-wise operations per parameter, i.e.\ $O(P)$, which is negligible compared to $3F$. Therefore
\[
\text{FLOPs per step} = F + 2F + O(P) \approx 3F.
\]
\end{answer}

\textbf{(d)} At 50\% MFU on an H100 (495\,TFLOP/s fp32), how many hours to train GPT-2 XL for 400K steps at batch 1024? (Backward = 2x forward FLOPs.) \\
\emph{Deliverable: The number of hours training would take, with a brief justification.}
\begin{answer}
\begin{align*}
\text{Effective throughput} &= 495 \times 0.5 = 247.5 \ \text{TFLOP/s} \\
\text{Forward per step (batch 1024)} &= 1024 \times F = 1024 \times 3.52 = 3604.48 \ \text{TFLOPs} \\
\text{FLOPs per step } (3F) &= 3 \times 3604.48 = 10813.44 \ \text{TFLOPs} \\
\text{Total FLOPs} &= 10813.44 \times 400{,}000 = 4.325 \times 10^{9} \ \text{TFLOPs} \\
\text{Time} &= \frac{4.325 \times 10^{9}}{247.5} \approx 1.748 \times 10^{7} \ \text{s} \approx \boxed{4854 \ \text{hours}}
\end{align*}
That is roughly $202$ days on a single H100. The forward pass scales with the batch of $1024$ sequences, the backward pass doubles it ($3F$ per step), over $400\text{K}$ steps, divided by the effective $247.5$ TFLOP/s throughput at $50\%$ MFU.
\end{answer}


\problem{learning\_rate\_schedule}{Implement cosine learning rate schedule with warmup (1 point)}
Write a function returning $\alpha_t$ given $t, \alpha_{\max}, \alpha_{\min}, T_w, T_c$.
\emph{Deliverable: Code that passes the provided tests.}
\begin{answer}
% TODO (optional notes)
\end{answer}


\problem{gradient\_clipping}{Implement gradient clipping (1 point)}
Implement in-place $\ell_2$-norm gradient clipping with $\varepsilon = 10^{-6}$.
\emph{Deliverable: Code that passes the provided tests.}
\begin{answer}
% TODO (optional notes)
\end{answer}


\section{Training Loop}

\problem{data\_loading}{Implement data loading (2 points)}
Write a function that samples a batch of (input, target) token-ID sequences from a NumPy array.
\emph{Deliverable: Code that passes the provided tests.}
\begin{answer}
% TODO (optional notes)
\end{answer}


\problem{checkpointing}{Implement model checkpointing (1 point)}
Implement \texttt{save\_checkpoint} and \texttt{load\_checkpoint}.
\emph{Deliverable: Code that passes the provided tests.}
\begin{answer}
% TODO (optional notes)
\end{answer}


\problem{training\_together}{Put it together (4 points)}
Write a training script supporting configurable hyperparameters, \texttt{np.memmap} data loading, checkpoint serialization, and periodic logging.
\emph{Deliverable: Training script.}
\begin{answer}
% TODO (optional notes)
\end{answer}


\section{Generating Text}

\problem{decoding}{Decoding (3 points)}
Implement a decoding function supporting prompt completion up to \texttt{<|endoftext|>} / max tokens, temperature scaling, and top-$p$ (nucleus) sampling.
\emph{Deliverable: Code.}
\begin{answer}
% TODO (optional notes)
\end{answer}


\section{Experiments}

\problem{experiment\_log}{Experiment logging (3 points)}
Build logging infrastructure and keep an experiment log for the problems in this section.
\emph{Deliverable: Logging infrastructure code and an experiment log document.}
\begin{answer}
Training runs are logged to Weights \& Biases (project \texttt{cs336-tinystories}). At every 100th step we log train loss, validation loss, and wall-clock time via \texttt{wandb.log(\{...\}, step=it)}, and the full hyperparameter config (model + optimizer dicts) is recorded at \texttt{wandb.init}. Runs execute on a cloud GPU through Modal (\texttt{modal\_train.py}), reading the pre-tokenized TinyStories arrays from a persistent volume. The main run trained a $d_{\text{model}}=512$, 4-layer model for 5000 steps, with validation loss decreasing from $\approx 9.2$ (random init, $\approx \ln 10000$) to $\approx 1.5$.
W\&B dashboard: \url{https://wandb.ai/rohan0401/cs336-tinystories/reports/Experiment-Log--VmlldzoxNzUyMDYxMw}
\end{answer}


\subsection*{7.2 TinyStories}

\problem{learning\_rate}{Tune the learning rate (3 points)}

\textbf{(a)} Perform a learning-rate sweep; report final losses / divergences. \\
\emph{Deliverable 1: Learning curves for multiple LRs. Explain your search strategy.} \\
\emph{Deliverable 2: A model with per-token validation loss $\leq 1.45$ on TinyStories (or $\leq 2.00$ on CPU/MPS).}
\begin{answer}
\textbf{Search strategy.} We swept $\text{lr}_{\max} \in \{10^{-4},\, 3\times10^{-4},\, 10^{-3},\, 3\times10^{-3},\, 10^{-2}\}$ (a geometric grid spanning two orders of magnitude), holding all else fixed (cosine schedule with 500 warmup steps, 5000 total steps, batch size 32), and compared validation-loss curves.

\textbf{Result.} Final validation losses: $\text{lr}=10^{-4}\!: 2.32$, $\ 3\times10^{-4}\!: 1.88$, $\ 10^{-3}\!: \mathbf{1.56}$, $\ 3\times10^{-3}\!: 1.86$, $\ 10^{-2}\!: 2.55$. The relationship is U-shaped: $\text{lr}_{\max}=10^{-3}$ gave the lowest loss and fastest stable convergence; smaller rates converged too slowly to reach a good loss within the step budget, and larger rates ($3\times10^{-3}$, $10^{-2}$) trained much worse as the steps began to overshoot. This is the expected pattern --- convergence speeds up with learning rate until a stability threshold, past which the loss climbs back up.

\begin{center}
\figph{figures/lr_sweep.png}
\end{center}
% export overlaid val-loss curves (per LR) from W&B as figures/lr_sweep.png
\end{answer}

\textbf{(b)} Investigate how the divergence point relates to the best LR (``edge of stability''). \\
\emph{Deliverable: Learning curves of increasing LR including at least one divergent run; analysis of the relation to convergence rates.}
\begin{answer}
The best-performing learning rate sits just below the point where training destabilizes, illustrating the ``edge of stability.'' In our sweep, $\text{lr}_{\max}=10^{-3}$ was the fastest rate that still trained cleanly; at $3\times10^{-3}$ and $10^{-2}$ the final loss rose again ($1.86$ and $2.55$) as the larger steps began to overshoot, and pushing the learning rate further (as seen with the largest batch/LR combination in the batch-size experiment) causes outright divergence. Rates well below the best value ($10^{-4}$) are stable but converge too slowly to reach a good loss within a fixed step budget. The practical implication is that the fastest useful learning rate is the largest one that does not degrade or destabilize training.
% overlaid curves showing 1e-3 (best) vs 1e-2 (much worse) --- see figures/lr_sweep.png
\end{answer}


\problem{batch\_size\_experiment}{Batch size variations (1 point)}
Vary batch size from 1 up to the GPU memory limit (include 64 and 128). \\
\emph{Deliverable 1: Learning curves for runs with different batch sizes (retune LR if needed).} \\
\emph{Deliverable 2: A few sentences discussing your findings.}
\begin{answer}
We swept batch size $\in \{1, 16, 64, 128, 256\}$ at a fixed budget of 5000 steps, scaling the learning rate with batch size ($\text{lr} = 10^{-3}\sqrt{B/32}$) since larger batches tolerate (and benefit from) a higher learning rate.

\begin{center}
\begin{tabular}{rrl}
\toprule
batch size & $\text{lr}_{\max}$ & final val loss \\
\midrule
1   & $1.8\times10^{-4}$ & $\approx 2.6$--$3.0$ \\
16  & $7.1\times10^{-4}$ & $1.91$ \\
64  & $1.4\times10^{-3}$ & $1.58$ \\
128 & $2.0\times10^{-3}$ & $\mathbf{1.48}$ \\
256 & $2.8\times10^{-3}$ & diverged (no valid loss logged) \\
\bottomrule
\end{tabular}
\end{center}

\textbf{Findings.} At a fixed step budget, larger batches reach lower validation loss (down to $1.48$ at $B=128$, vs.\ $\approx 3$ at $B=1$), for two reasons: a larger batch processes proportionally more total tokens in the same number of steps, and it gives lower-variance gradient estimates that support the larger, more stable learning rate. However, this is a fixed-\emph{steps} comparison, not fixed-\emph{compute}: each larger-batch step does more work, so part of the gain simply reflects more total computation. At the largest setting ($B=256$, $\text{lr}=2.8\times10^{-3}$) the LR scaling overshot and training diverged, illustrating that batch-size scaling must be paired with careful LR retuning.

\begin{center}
\figph{figures/batch_size.png}
\end{center}
% export the overlaid val-loss curves (per batch size) from W&B as figures/batch_size.png
\end{answer}


\problem{generate}{Generate text (1 point)}
Report text generated by your trained model (tune temperature / top-$p$ for fluency). \\
\emph{Deliverable: $\geq 256$ tokens of generated text (or until first \texttt{<|endoftext|>}), plus a brief comment on fluency and at least two factors affecting quality.}
\begin{answer}
Generated from the trained TinyStories model with the prompt \texttt{"Once upon a time"}, temperature $0.8$, top-$p$ $0.9$, sampling until the first \texttt{<|endoftext|>}:

\begin{lstlisting}
Once upon a time, there was a beautiful fairy. She lived in a big,
pretty garden. The fairy had a pretty garden with many flowers. She
liked to look at the flowers and the bugs. One day, the fairy met a
boy. The boy was deaf. He could not hear her. The fairy wanted to
help the boy. She said, "Please let me play with you. I will teach
you how to be nice." The boy was happy. He started to play with the
fairy. They played and had fun. The fairy and the boy became good
friends. They played together every day. And they lived happily
ever after.
\end{lstlisting}

\textbf{Fluency.} The output is fluent and coherent: grammatical sentences, correctly punctuated dialogue, a consistent cast, and a complete narrative arc (setup $\to$ problem $\to$ resolution $\to$ ending), terminating naturally on \texttt{<|endoftext|>}. The simple vocabulary and story structure are exactly on-distribution for TinyStories.

\textbf{Factors affecting quality.} (1) \emph{Decoding temperature / top-$p$:} at low temperature ($0.1$) generation collapses toward greedy and becomes repetitive; at high temperature ($1.0$) it is more creative but less coherent (occasional non-sequiturs); $\tau=0.8$ with top-$p=0.9$ balances diversity and fluency. (2) \emph{Model capacity and training:} the model reached validation loss $\approx 1.5$; a larger $d_{\text{model}}$/more layers or more training steps would sharpen next-token predictions and further improve fluency. (3) \emph{Training-data domain:} because it was trained only on TinyStories, it produces simple children's-story prose and would not generalize to other domains.
\end{answer}


\subsection*{7.3 Ablations and Architecture Modification}

\problem{layer\_norm\_ablation}{Remove RMSNorm and train (1 point)}
Remove all RMSNorms; train at the previous optimal LR; try a lower LR. \\
\emph{Deliverable 1: Learning curves (at prior optimal LR and at best LR).} \\
\emph{Deliverable 2: A few sentences of commentary on the impact of RMSNorm.}
\begin{answer}
Removing all RMSNorms (both per-block norms and the final norm, replaced by identity) degrades the final validation loss from the baseline $1.564$ to $\mathbf{1.715}$ ($+0.15$). Notably, at this small scale the no-norm model still trained stably at the baseline learning rate ($10^{-3}$) rather than diverging, so no lower learning rate was required to keep it stable; normalization here improves convergence quality rather than being strictly necessary for stability. RMSNorm helps by keeping activation magnitudes controlled across depth, which improves gradient flow and lets the model reach a lower loss within the same step budget.

\begin{center}
\figph{figures/ablation_rmsnorm.png}
\end{center}
% export the overlaid curves (baseline vs. no-RMSNorm) from W&B as figures/ablation_rmsnorm.png
\end{answer}


\problem{pre\_norm\_ablation}{Implement post-norm and train (1 point)}
Convert to post-norm; train. \\
\emph{Deliverable: A learning curve for post-norm vs. pre-norm.}
\begin{answer}
Switching from pre-norm ($y = x + \text{sublayer}(\text{norm}(x))$) to post-norm ($y = \text{norm}(x + \text{sublayer}(x))$) raises the final validation loss from $1.564$ to $\mathbf{1.769}$ ($+0.21$). Pre-norm is more effective because the residual stream stays un-normalized, giving a clean identity path for gradients to flow through the depth of the network; in post-norm the normalization sits on the residual path, which attenuates that signal and slows convergence. This matches why modern LLMs (GPT-3, LLaMA, PaLM) standardized on pre-norm.

\begin{center}
\figph{figures/ablation_prenorm.png}
\end{center}
% export the overlaid curves (pre-norm vs. post-norm) from W&B as figures/ablation_prenorm.png
\end{answer}


\problem{no\_pos\_emb}{Implement NoPE (1 point)}
Remove RoPE (no position embeddings at all). \\
\emph{Deliverable: A learning curve comparing RoPE vs. NoPE.}
\begin{answer}
Removing all positional information (NoPE --- no RoPE) raises the final validation loss from $1.564$ to $\mathbf{1.802}$ ($+0.24$), the largest degradation of any single-component ablation. This is expected: self-attention is permutation-invariant, so without a positional signal the model cannot distinguish token order (e.g.\ ``dog bites man'' vs.\ ``man bites dog''). Interestingly the model does not collapse entirely --- because of \emph{causal masking}, each position attends to a different-sized prefix, which leaks a weak positional signal, so a decoder-only LM degrades but remains partially functional without explicit position embeddings.

\begin{center}
\figph{figures/ablation_nope.png}
\end{center}
% export the overlaid curves (RoPE vs. NoPE) from W&B as figures/ablation_nope.png
\end{answer}


\problem{swiglu\_ablation}{SwiGLU vs. SiLU (1 point)}
Replace SwiGLU with a SiLU FFN (no GLU), setting $d_{\text{ff}} = 4 d_{\text{model}}$ to roughly match parameter count. \\
\emph{Deliverable 1: A learning curve comparing SwiGLU vs. SiLU.} \\
\emph{Deliverable 2: A few sentences discussing your findings.}
\begin{answer}
Replacing SwiGLU (gated, $d_{\text{ff}}=1344$, three matrices) with a plain SiLU feed-forward (no gate, $d_{\text{ff}}=2048$, two matrices) raises the final validation loss from $1.564$ to $\mathbf{1.642}$ ($+0.08$). The two variants have nearly identical parameter counts ($\approx 2.06$M vs.\ $2.10$M per block), so the gap is attributable to the \textbf{GLU gating mechanism itself}, not to parameter count. This matches Shazeer's finding that gating improves language-modeling performance; it is the smallest of our ablation effects, indicating the gate helps but is less critical than positional information or normalization at this scale.

\begin{center}
\figph{figures/ablation_swiglu.png}
\end{center}
% export the overlaid curves (SwiGLU vs. SiLU) from W&B as figures/ablation_swiglu.png
\end{answer}


\subsection*{7.4 Running on OpenWebText}

\problem{main\_experiment}{Experiment on OWT (2 points)}
Train on OpenWebText with the same architecture / total iterations as TinyStories. \\
\emph{Deliverable 1: A learning curve on OWT; describe the loss difference vs. TinyStories --- how to interpret?} \\
\emph{Deliverable 2: Generated text from the OWT LM; comment on fluency and why quality differs despite equal compute.}
\begin{answer}
Training the same architecture (4 layers, $d_{\text{model}}=512$) for the same 5000 steps on OpenWebText (32k vocabulary) reaches a final validation loss of $\approx \mathbf{4.58}$, far higher than the $\approx 1.5$ obtained on TinyStories.

\textbf{Interpretation.} The gap is expected despite equal compute. TinyStories is a narrow, highly regular corpus (simple children's stories with a small vocabulary and repetitive structure), so a small model can fit it tightly. OpenWebText is scraped web text: vastly more diverse in topic, style, language, and formatting, with a larger 32k vocabulary. The same-capacity model therefore cannot compress OWT nearly as well, and even the random-init baseline is higher ($\ln 32000 \approx 10.4$ vs.\ $\ln 10000 \approx 9.2$). In short, loss measures how predictable the data is under the model, and OWT is intrinsically far less predictable than TinyStories.

\textbf{Generated text.} Text sampled from the OWT model is noticeably less fluent than the TinyStories model at equal compute --- it produces locally plausible but globally incoherent web-style fragments --- because the model has spread its limited capacity over a much broader and less repetitive distribution, so it has not learned any single domain well.

\begin{center}
\figph{figures/owt_curve.png}
\end{center}
% export the OWT vs. TinyStories val-loss curves from W&B as figures/owt_curve.png;
% paste an OWT generation sample here in an lstlisting block
\end{answer}


\subsection*{7.5 Your Own Modification + Leaderboard}

\problem{leaderboard}{Leaderboard (6 points)}
Train a model under leaderboard rules (B200, $\leq 45$ minutes, OWT only) aiming to minimize validation loss within 0.75 B200-hours. Target: beat the 5.0 naive baseline. \\
\emph{Deliverable: Final validation loss, a learning curve with wall-clock-time x-axis $< 45$ minutes, and a description of what you did.}
\begin{answer}
\textbf{Final validation loss:} $\mathbf{4.48}$ on OpenWebText (best $4.45$ at step 4500), beating the naive $5.0$ baseline. The full 5000-step run was on a single \textbf{NVIDIA B200} and took \textbf{2.5 minutes} of wall-clock time ($151$\,s), well within the 45-minute / 0.75 B200-hour budget.

\textbf{What we did.} We trained the pre-norm Transformer LM (4 layers, $d_{\text{model}}=512$, 16 heads, $d_{\text{ff}}=1344$, SwiGLU, RoPE, RMSNorm; OWT 32k vocabulary) with AdamW, a cosine learning-rate schedule with 500 warmup steps, gradient clipping at $\ell_2$-norm $1.0$, and $\text{lr}_{\max}=10^{-3}$ (the value that won the TinyStories sweep). Two throughput optimizations kept per-step time low on the B200: \texttt{torch.compile} (kernel fusion via the inductor backend) and mixed-precision \texttt{bfloat16} autocast on the forward pass. These do not change the loss trajectory --- bf16 has the same exponent range as fp32, and our RMSNorm upcasts internally for numerical stability. Data was streamed from a memory-mapped \texttt{uint16} token array, so the 5.1\,GB OWT training set never had to fit in memory.

\textbf{Learning curve.} Over the full 5000-step run, validation loss fell from the random-init value ($\approx \ln 32000 \approx 10.4$) to $4.48$; the curve below uses a wall-clock-time x-axis confirming the run finished in $\approx 2.5$ minutes.

\begin{center}
\figph{figures/leaderboard_curve.png}
\end{center}
% figures/leaderboard_curve.png: export the val-loss-vs-wall_clock_s curve from W&B run 21dar844
\end{answer}


\end{document}
